\usepackage{cmap}
\usepackage[T2A]{fontenc}
\usepackage[utf8]{inputenc}
\usepackage[english,russian]{babel}



\usepackage[
    backend=biber,
    style=gost-numeric,
    sorting=none,
    language=russian,
    url=true,
    doi=true,
    date=year,
    maxnames=3,
    minnames=1,
    giveninits=true
]{biblatex}

\addbibresource{references.bib}

% Русские строки
\DefineBibliographyStrings{russian}{
  andothers = {и др.},
  urlseen   = {дата обращения\addcolon},
}

% Формат имён
\DeclareNameAlias{sortname}{family-given}
\DeclareNameAlias{default}{family-given}

% Убираем лишние поля
\AtEveryBibitem{
  \clearfield{isbn}
  \clearfield{issn}
  \clearfield{abstract}
  \clearfield{series}
}

% Кастомный заголовок
\defbibheading{gostbib}[\bibname]{%
  \chapter*{#1}%
  \addcontentsline{toc}{chapter}{#1}%
  \markboth{#1}{}%
}
% === ИСПРАВЛЕНИЕ ПЕРЕНОСОВ В БИБЛИОГРАФИИ ===

% 1. Делаем шрифт URL таким же, как основной текст (чтобы переносы работали корректно)
\urlstyle{same}

% 2. Разрешаем переносы ВНУТРИ URL по символам: /, -, _, ., :, ~, ?
% Это предотвращает вылезание ссылок за правое поле
\makeatletter
\def\UrlBreaks{
  \do\/\do-\do_\do.\do:\do\@\do\\\do\/\do\%\do\~\do\?\do\&\do\=\do\#\do\+\do\*\do\,\do\;\do\|\do\<\do\>\do\^\do\'\do\"\do\`\do\!\do\$\do\[\do\]\do\{\do\}\do\(\do\)
}
\makeatother
% === Формат ссылок: Режим доступа: URL: ... ===
\DeclareFieldFormat[online,misc]{url}{%
  Режим доступа: URL\addcolon\space\url{#1}%
}
\DeclareFieldFormat[inproceedings,misc]{url}{%
  Режим доступа: URL\addcolon\space\url{#1}%
}


% Разделители
\renewcommand*{\newunitpunct}{\addperiod\space}
\renewcommand*{\finentrypunct}{\addperiod}

\renewcommand*{\mkbibnamefamily}[1]{#1}

\usepackage{hyperref}
\hypersetup{
    hidelinks,
    unicode,
    pdfencoding=auto
}

\usepackage{multirow}
\usepackage{threeparttable}

\usepackage[14pt]{extsizes}
\usepackage{caption}
\captionsetup[table]{justification=raggedright,singlelinecheck=off}
\captionsetup{labelsep=endash}
\captionsetup[figure]{name={Рисунок}}

\usepackage{amsmath}
\usepackage{enumitem} 
\setenumerate[0]{label=\arabic*)}
\renewcommand{\labelitemi}{---}

\usepackage[%
    a4paper,
    left=30mm,      % поле под переплёт (внутреннее при twoside)
    right=15mm,     % внешнее поле
    top=20mm,
    bottom=20mm,
    bindingoffset=0mm,  % дополнительный отступ на переплёт (обычно 0, если left=30 уже учтено)
    oneside,        % зеркальные поля для двусторонней печати
    % showframe,    % РАСКОММЕНТИРУЙТЕ чтобы увидеть границы полей при отладке. Потом закомментируйте!
]{geometry}

\usepackage{float}

\makeatletter
\renewcommand\LARGE{\@setfontsize\LARGE{22pt}{20}}
\renewcommand\Large{\@setfontsize\Large{20pt}{20}}
\renewcommand\large{\@setfontsize\large{16pt}{20}}
\makeatother

\usepackage{titlesec}

\titleformat{\chapter}[block]
  {\hspace{\parindent}\normalsize\bfseries}   % Шрифт номера
  {\thechapter}
  {1em}                                       % Отступ после номера
  {\normalsize\bfseries\raggedright}          % Шрифт заголовка (ТЕПЕРЬ 14pt)

% === ГЛАВЫ БЕЗ НОМЕРА (Введение, Заключение, Список литературы и т.д.) ===
\titleformat{name=\chapter,numberless}[block]
  {\centering\normalsize\bfseries}
  {}
  {0pt}
  {}

% === РАЗДЕЛЫ ===
\titleformat{\section}[block]
  {\hspace{\parindent}\normalsize\bfseries}{\thesection}{1em}{\normalsize\bfseries\raggedright}

% === ПОДРАЗДЕЛЫ ===
\titleformat{\subsection}[block]
  {\hspace{\parindent}\normalsize\bfseries}{\thesubsection}{1em}{\normalsize\bfseries\raggedright}

% === ВЕРТИКАЛЬНЫЕ ОТСТУПЫ ===
\titlespacing*{\chapter}{12.5mm}{0pt}{10pt}
\titlespacing*{\section}{12.5mm}{10pt}{10pt}
\titlespacing*{\subsection}{12.5mm}{10pt}{10pt}
\titlespacing*{\subsubsection}{12.5mm}{10pt}{10pt}

\makeatother
\frenchspacing
\usepackage{indentfirst}
\setlength\parindent{1.25cm}

\usepackage[figure,table]{totalcount}
\usepackage{rotating}
\usepackage{lastpage}

\usepackage{array}
\newenvironment{signstabular}[1][1]{
	\renewcommand*{\arraystretch}{#1}
	\tabular
}{
	\endtabular
}

\usepackage{csvsimple}
\usepackage{xcolor}

\usepackage{pgfplots}
\pgfplotsset{compat=1.18}
\usetikzlibrary{datavisualization}
\usetikzlibrary{datavisualization.formats.functions}

\usepackage{graphicx}
\newcommand{\imgScale}[3] {
	\begin{figure}[h!]
		\center{\includegraphics[scale=#1]{img/#2}}
		\caption{#3}
		\label{img:#2}
	\end{figure}
}

\newcommand{\imgHeight}[3] {
	\begin{figure}[h!]
		\center{\includegraphics[height=#1]{img/#2}}
		\caption{#3}
		\label{img:#2}
	\end{figure}
}

\usepackage[justification=centering]{caption}
\hypersetup{hidelinks}
\newcommand{\code}[1]{\texttt{#1}}

\usepackage{tocloft}

% === ЗАГОЛОВОК ===
\renewcommand{\contentsname}{СОДЕРЖАНИЕ}
\renewcommand{\cfttoctitlefont}{\hspace*{\fill}\large\bfseries} % Центрирование, размер как у заголовков
\renewcommand{\cftaftertoctitle}{\hspace*{\fill}}
\setlength{\cftbeforetoctitleskip}{0pt}
\setlength{\cftaftertoctitleskip}{10pt}

% === ШРИФТ ЭЛЕМЕНТОВ ===
\renewcommand{\cftchapfont}{\large\bfseries}       % Главы
\renewcommand{\cftsecfont}{\large}                 % Разделы
\renewcommand{\cftsubsecfont}{\large}              % Подразделы
\renewcommand{\cftchappagefont}{\large\bfseries}   % Номера страниц глав
\renewcommand{\cftsecpagefont}{\large}             % Номера страниц разделов
\renewcommand{\cftsubsecpagefont}{\large}          % Номера страниц подразделов

% === ОТСТУПЫ СЛЕВА (кратность 1.25 см) ===
\setlength{\cftchapindent}{0pt}
\setlength{\cftsecindent}{1.25cm}
\setlength{\cftsubsecindent}{2.5cm}
\setlength{\cftsubsubsecindent}{3.75cm}

% === ТОЧКИ И НОМЕРА ===
\renewcommand{\cftdotsep}{4.5}                     % Расстояние между точками
\setlength{\cftchapnumwidth}{1.5cm}                % Фиксированная ширина под номер главы
\setlength{\cftsecnumwidth}{1.5cm}
\setlength{\cftsubsecnumwidth}{1.5cm}

% === ВЕРТИКАЛЬНЫЕ ИНТЕРВАЛЫ ===
\setlength{\cftbeforechapskip}{0pt}
\setlength{\cftbeforesecskip}{0pt}
\setlength{\cftbeforesubsecskip}{0pt}

% === ФОРМАТ НОМЕРА ГЛАВЫ ===
\renewcommand{\cftchappresnum}{}    % Убираем слово "Глава"
\renewcommand{\cftchapaftersnum}{} % Точка после номера

\renewcommand{\cftchapleader}{\cftdotfill{\cftdotsep}}
\usepackage{csquotes}
\usepackage{setspace}
\onehalfspacing
\usepackage{makecell}
\renewcommand\theadfont{\bfseries\small}

\usepackage{listings}
\usepackage{xcolor}

\lstset{
    language=Python,
    inputencoding=utf8,
    extendedchars=true,
    basicstyle=\ttfamily\small,
    keywordstyle=\color{blue}\bfseries,
    commentstyle=\color{gray}\itshape,
    stringstyle=\color{orange},
    numbers=left,
    numberstyle=\tiny\color{black},
    stepnumber=1,
    numbersep=5pt,
    % === РАМКА И ПЕРЕНОСЫ ===
    frame=single,               % Полная рамка со всех сторон
    framerule=0.4pt,            % Толщина линии рамки
    framesep=3pt,               % Отступ текста от рамки
    rulecolor=\color{black},    % Цвет рамки
    breaklines=true,            % Автоматический перенос длинных строк
    breakatwhitespace=true,     % Перенос только по пробелам
    % === ПОДПИСИ ===
    captionpos=t,               % Подпись СВЕРХУ (ГОСТ)
    % === ПРОЧЕЕ ===
    showspaces=false,
    showstringspaces=false,
    showtabs=false,
    tabsize=4,
    literate=
        {а}{{\selectfont\char224}}1 {б}{{\selectfont\char225}}1 {в}{{\selectfont\char226}}1
        {г}{{\selectfont\char227}}1 {д}{{\selectfont\char228}}1 {е}{{\selectfont\char229}}1
        {ё}{{\selectfont\char184}}1 {ж}{{\selectfont\char230}}1 {з}{{\selectfont\char231}}1
        {и}{{\selectfont\char232}}1 {й}{{\selectfont\char233}}1 {к}{{\selectfont\char234}}1
        {л}{{\selectfont\char235}}1 {м}{{\selectfont\char236}}1 {н}{{\selectfont\char237}}1
        {о}{{\selectfont\char238}}1 {п}{{\selectfont\char239}}1 {р}{{\selectfont\char240}}1
        {с}{{\selectfont\char241}}1 {т}{{\selectfont\char242}}1 {у}{{\selectfont\char243}}1
        {ф}{{\selectfont\char244}}1 {х}{{\selectfont\char245}}1 {ц}{{\selectfont\char246}}1
        {ч}{{\selectfont\char247}}1 {ш}{{\selectfont\char248}}1 {щ}{{\selectfont\char249}}1
        {ъ}{{\selectfont\char250}}1 {ы}{{\selectfont\char251}}1 {ь}{{\selectfont\char252}}1
        {э}{{\selectfont\char253}}1 {ю}{{\selectfont\char254}}1 {я}{{\selectfont\char255}}1
        {А}{{\selectfont\char192}}1 {Б}{{\selectfont\char193}}1 {В}{{\selectfont\char194}}1
        {Г}{{\selectfont\char195}}1 {Д}{{\selectfont\char196}}1 {Е}{{\selectfont\char197}}1
        {Ё}{{\selectfont\char168}}1 {Ж}{{\selectfont\char198}}1 {З}{{\selectfont\char199}}1
        {И}{{\selectfont\char200}}1 {Й}{{\selectfont\char201}}1 {К}{{\selectfont\char202}}1
        {Л}{{\selectfont\char203}}1 {М}{{\selectfont\char204}}1 {Н}{{\selectfont\char205}}1
        {О}{{\selectfont\char206}}1 {П}{{\selectfont\char207}}1 {Р}{{\selectfont\char208}}1
        {С}{{\selectfont\char209}}1 {Т}{{\selectfont\char210}}1 {У}{{\selectfont\char211}}1
        {Ф}{{\selectfont\char212}}1 {Х}{{\selectfont\char213}}1 {Ц}{{\selectfont\char214}}1
        {Ч}{{\selectfont\char215}}1 {Ш}{{\selectfont\char216}}1 {Щ}{{\selectfont\char217}}1
        {Ъ}{{\selectfont\char218}}1 {Ы}{{\selectfont\char219}}1 {Ь}{{\selectfont\char220}}1
        {Э}{{\selectfont\char221}}1 {Ю}{{\selectfont\char222}}1 {Я}{{\selectfont\char223}}1,
}
\makeatletter
\newcommand{\lstcontinuation}{%
  \addtocounter{lstlisting}{-1}% Откатываем счётчик, чтобы номер остался A.1
  \renewcommand{\lstlistingname}{Продолжение листинга}% Меняем префикс
}
\newcommand{\lstendcontinuation}{%
  \renewcommand{\lstlistingname}{Листинг}% Возвращаем слово "Листинг"
  \lstset{labelsep=endash}% Возвращаем тире для следующих листингов
}
\makeatother
\usepackage{caption}
\captionsetup[lstlisting]{font=normalsize, labelfont=bf, labelsep=endash}

\usepackage{amsmath,amssymb,amsfonts}
\usepackage{icomma}
\usepackage{upgreek}

\newcommand{\imgfig}[1]{
    \includegraphics[
        width=\textwidth,
        height=0.85\textheight,
        keepaspectratio
    ]{#1}
}
\newcommand{\fixunderline}[3]
{
	$\underset{\text{#3}}{\text{\stackengine{0pt}{\hspace{#2}}{\text{#1}}{O}{c}{F}{F}{L}}}$
}

% Команда для поля с подписью/датой ПОД подчёркиванием
% Добавлен \vphantom для выравнивания по высоте с полем фамилии
\newcommand{\signaturefield}{%
	\renewcommand{\arraystretch}{1}%
	\begin{tabular}[t]{@{}c@{}}%
		\vphantom{Фамилия~И.~О.}\\[-0.2em]%
		\underline{\hspace{4cm}}\\[-0.2em]%
		\scriptsize (Подпись, дата)%
	\end{tabular}%
}

% Команда для поля с фамилией НАД подчёркиванием
\newcommand{\signfield}[1]{%
	\renewcommand{\arraystretch}{1}%
	\begin{tabular}[t]{@{}c@{}}%
		#1\\[-0.2em]%
		\underline{\hspace{4cm}}%
	\end{tabular}%
}

% Команда для пустого поля (только линия)
\newcommand{\emptysignfield}{%
	\renewcommand{\arraystretch}{1}%
	\begin{tabular}[t]{@{}c@{}}%
		\vphantom{Фамилия~И.~О.}\\[-0.2em]%
		\underline{\hspace{4cm}}%
	\end{tabular}%
}